\subsection{Gwarancje probabilistyczne i błąd aproksymacji – Sformułowanie i dowody głównych twierdzeń ograniczających błąd aproksymacji $\|f - \hat{f}\|_\infty$ z odpowiednio wysokim prawdopodobieństwem.}

W niniejszym podrozdziale przedstawiamy rygorystyczne gwarancje sukcesu opisanych wcześniej algorytmów. Ze względu na losowy dobór punktów próbkowania i kierunków pomiarowych, błąd aproksymacji szacowany jest z wysokim prawdopodobieństwem. Analizę przeprowadzimy osobno dla modelu jednowymiarowego ($k=1$) oraz wielowymiarowego ($k \ge 1$).Przypadek jednowymiarowy ($k=1$) Zacznijmy od sformułowania i udowodnienia głównego twierdzenia dotyczącego Algorytmu 1.

\begin{tw}
     Niech $0 < s < 1$ oraz $\log d \le m_\Phi \le [\log 6]^{-2}d$. Wówczas istnieje stała $c_1'$, taka że przy użyciu $m_\mathcal{X} \cdot (m_\Phi + 1)$ ewaluacji funkcji $f$, Algorytm 1 definiuje funkcję $\hat{f}: B_{\mathbb{R}^d}(1+\bar{\epsilon}) \to \mathbb{R}$, która z prawdopodobieństwem wynoszącym co najmniej:
$$1 - \left(e^{-c_1' m_\Phi} + e^{-\sqrt{m_\Phi d}} + 2e^{-\frac{2m_\mathcal{X} s^2 \alpha^2}{C_2^4}}\right)$$spełnia ograniczenie błędu:
$$\|f - \hat{f}\|_\infty \le 2C_2(1+\bar{\epsilon})\frac{\nu_1}{\sqrt{\alpha(1-s)} - \nu_1}$$gdzie parametr $\nu_1$ jest zdefiniowany jako:
$$\nu_1 = C' \left( \left[ \frac{m_\Phi}{\log(d/m_\Phi)} \right]^{1/2 - 1/q} + \frac{\epsilon}{\sqrt{m_\Phi}} \right)$$a $C'$ zależy wyłącznie od stałych $C_1$ i $C_2$ (ograniczających odpowiednio normę wektora i pochodne funkcji $g$).
\end{tw}

\begin{proof}
Dowód opiera się na połączeniu trzech kluczowych oszacowań: stabilności próbkowania skompresowanego, Nierówności Hoeffdinga dla próbkowania sferycznego oraz stabilności podprzestrzeni.
    
W pierwszej kolejności musimy zagwarantować, że maksimum z pochodnych, czyli $\max_{j=1,\dots,m_\mathcal{X}} |g'(a \cdot \xi_j)|$, jest odpowiednio duże, co uchroni nas przed szumem. Ponieważ zmienne losowe $X_j = |g'(a \cdot \xi_j)|^2$ są ograniczone do przedziału $[0, C_2^2]$ i ich wartość oczekiwana to $\alpha > 0$, możemy zastosować klasyczną Nierówność Hoeffdinga. Na jej mocy otrzymujemy, że:
$$\max_{j=1,\dots,m_\mathcal{X}} |g'(a \cdot \xi_j)| \ge \sqrt{\alpha(1-s)}$$z prawdopodobieństwem $1 - 2e^{-\frac{2m_\mathcal{X} s^2 \alpha^2}{C_2^4}}$.
Wiemy z własności ograniczonej izometrii (RIP) dla macierzy losowych (Wniosek 3.3 w pracy bazowej), że błąd rekonstrukcji wektora rzadkiego przy pomocy minimalizacji $l_1$ jest stabilny. Oznacza to, że z prawdopodobieństwem co najmniej $1 - (e^{-c_1' m_\Phi} + e^{-\sqrt{m_\Phi d}})$ dla odzyskanych wektorów zachodzi:
$$\|x_j - \hat{x}_j\|_{l_2^d} \le \nu_1$$Krok 3: Łącząc powyższe z Lematem o stabilności podprzestrzeni, dla znormalizowanego wektora $\hat{a}$ wybranego w algorytmie zachodzi oszacowanie odległości od prawdziwego wektora $a$:
$$\|\text{sign}(g'(a \cdot \xi_{j_0}))\hat{a} - a\|_{l_2^d} \le \frac{2\nu_1}{\sqrt{\alpha(1-s)} - \nu_1}$$z prawdopodobieństwem łącznym $1 - (e^{-c_1' m_\Phi} + e^{-\sqrt{m_\Phi d}} + 2e^{-\frac{2m_\mathcal{X} s^2 \alpha^2}{C_2^4}})$. Znak jest wynikiem niejednoznaczności kierunku wektora.Krok 4: Ostatecznie szacujemy błąd jednostajny funkcji. Dla dowolnego $x \in B_{\mathbb{R}^d}(1+\bar{\epsilon})$:
$$|f(x) - \hat{f}(x)| = |g(a \cdot x) - \hat{g}(\hat{a} \cdot x)| = |g(a \cdot x) - f(\hat{a}^T \hat{a} x)| = |g(a \cdot x) - g(a \hat{a}^T \hat{a} x)|$$Korzystając z Twierdzenia o wartości średniej i faktu, że $\hat{a}(I_d - \hat{a}^T \hat{a}) = 0$, wprowadzamy składową ze znakiem, uzyskując:
$$|f(x) - \hat{f}(x)| \le C_2 |a \cdot (I_d - \hat{a}^T \hat{a})x| = C_2 |(a - \text{sign}(g'(a \cdot \xi_{j_0}))\hat{a}) \cdot (I_d - \hat{a}^T \hat{a})x|$$Stosując nierówność Cauchy'ego-Schwarza i wprowadzając oszacowanie z Kroku 3, otrzymujemy finalny błąd:
$$|f(x) - \hat{f}(x)| \le C_2 \|a - \text{sign}(g'(a \cdot \xi_{j_0}))\hat{a}\|_{l_2^d} \cdot \|x\|_{l_2^d} \le 2C_2(1+\bar{\epsilon})\frac{\nu_1}{\sqrt{\alpha(1-s)} - \nu_1}$$Co kończy dowód. $\blacksquare$
\end{proof}
Przypadek uogólniony ($k \ge 1$)Analogiczne twierdzenie formułujemy dla Algorytmu 2, z tym że wymaga ono zaawansowanych macierzowych nierówności probabilistycznych.Twierdzenie 4.1. Niech $\log d \le m_\Phi \le [\log 6]^2 d$. Wówczas istnieje stała $c_1'$, taka że wykorzystując $m_\mathcal{X} \cdot (m_\Phi + 1)$ ewaluacji funkcji $f$, Algorytm 2 definiuje funkcję $\hat{f}$, która z prawdopodobieństwem:
$$1 - \left(e^{-c_1' m_\Phi} + e^{-\sqrt{m_\Phi d}} + k e^{-\frac{m_\mathcal{X} \alpha s^2}{2kC_2^2}}\right)$$spełnia:
$$\|f - \hat{f}\|_\infty \le 2C_2\sqrt{k}(1+\bar{\epsilon})\frac{\nu_2}{\sqrt{\alpha(1-s)} - \nu_2}$$gdzie:
$$\nu_2 = C\left(k^{1/q}\left[\frac{m_\Phi}{\log(d/m_\Phi)}\right]^{1/2 - 1/q} + \frac{\epsilon k^2}{\sqrt{m_\Phi}}\right)$$Dowód Twierdzenia 4.1.
W tym przypadku centralnym problemem nie jest wybór jednego wektora, lecz stabilność całej podprzestrzeni rozpiętej przez wektory osobliwe.Krok 1: Wpierw szacujemy normę Frobeniusa błędu całej macierzy $X$. Stosując standardowe wyniki dla próbkowania skompresowanego kolumna po kolumnie, otrzymujemy:
$$\|X - \hat{X}\|_F \le \sqrt{m_\mathcal{X}}\nu_2$$z prawdopodobieństwem $1 - (e^{-c_1' m_\Phi} + e^{-\sqrt{m_\Phi d}})$.Krok 2: Musimy upewnić się, że najmniejsza niezerowa wartość osobliwa macierzy $X$ nie "zapada się" w kierunku zera. Macierz $X^T X$ jest sumą niezależnych, losowych macierzy półokreślonych dodatnio. Korzystając z macierzowej nierówności Chernoffa w ujęciu Troppa, udowadnia się (Lemat 4.5), że dla $\sigma_k(X^T)$ zachodzi:
$$\sigma_k(X^T) \ge \sqrt{m_\mathcal{X}\alpha(1-s)}$$z prawdopodobieństwem $1 - k e^{-\frac{m_\mathcal{X} \alpha s^2}{2kC_2^2}}$.Krok 3: Stabilność podprzestrzeni gwarantuje Twierdzenie Wedina o ograniczeniu perturbacji. Mówi ono, że dla odzyskanej bazy prawych wektorów osobliwych $\hat{V}_1$ oraz prawdziwej bazy $V_1$ mamy:
$$\|V_1 V_1^T - \hat{V}_1 \hat{V}_1^T\|_F \le \frac{2\sqrt{m_\mathcal{X}}\nu_2}{\sigma_k(X^T) - \sqrt{m_\mathcal{X}}\nu_2}$$Wstawiając ograniczenie z Kroku 2, redukujemy czynnik $\sqrt{m_\mathcal{X}}$ uzyskując:
$$\|V_1 V_1^T - \hat{V}_1 \hat{V}_1^T\|_F \le \frac{2\nu_2}{\sqrt{\alpha(1-s)} - \nu_2}$$Co oznacza, że również $\|\hat{A}^T \hat{A} - A^T A\|_F$ jest ograniczone tą samą wielkością (z łącznym prawdopodobieństwem podanym w Tezie).Krok 4: Aby zakończyć, przeprowadzamy łańcuch oszacowań błędu funkcji, biorąc pod uwagę, że dla ortonormalnej macierzy zachodzi $A = A A^T A$:
$$|f(x) - \hat{f}(x)| = |g(Ax) - g(A\hat{A}^T \hat{A}x)| \le C_2\sqrt{k} \|Ax - A\hat{A}^T\hat{A}x\|_{l_2^k}$$Przekształcając dalej z własności norm macierzowych:
$$|f(x) - \hat{f}(x)| \le C_2\sqrt{k} \|A(A^T A - \hat{A}^T \hat{A})x\|_{l_2^k} \le C_2\sqrt{k} \|A^T A - \hat{A}^T \hat{A}\|_F \|x\|_{l_2^d}$$Wstawiając ostatecznie ograniczenie na $x \in B_{\mathbb{R}^d}(1+\bar{\epsilon})$ oraz wynik z Twierdzenia Wedina, otrzymujemy tezę:
$$|f(x) - \hat{f}(x)| \le 2C_2\sqrt{k}(1+\bar{\epsilon})\frac{\nu_2}{\sqrt{\alpha(1-s)} - \nu_2}$$Co kończy dowód Twierdzenia 4.1. $\blacksquare$